\documentclass[12pt]{amsart}
\pdfoutput=1
\usepackage{latexsym, amsmath, amssymb, amsthm, bm}
\usepackage{calrsfs}
\usepackage{subfig}

\usepackage[mathscr]{eucal}
\usepackage{mathtools}
\usepackage{paralist}
\usepackage[alphabetic]{amsrefs}
\usepackage[T1]{fontenc}
\usepackage{mathptmx}
\usepackage{microtype}
\usepackage{tikz}
\renewcommand{\baselinestretch}{1.05}
\usepackage[colorlinks=true,linkcolor=blue,urlcolor=blue, citecolor=blue]%
    {hyperref}
\linespread{1.06}
%%--LAYOUT--------------------------------------------------------------------
\usepackage[paperwidth=8.5in, paperheight=11in, inner=2.5cm, outer=2.5cm,% 
    top=2.5cm, bottom=2.5cm]{geometry}
%%--OTHER ENVIRONMENTS--------------------------------------------------------
\newtheorem{lemma}{Lemma}[section]
\newtheorem{theorem}[lemma]{Theorem}
\newtheorem{corollary}[lemma]{Corollary}
\newtheorem{proposition}[lemma]{Proposition}
\theoremstyle{definition}
\newtheorem{definition}[lemma]{Definition}
\newtheorem{remark}[lemma]{Remark}
\newtheorem{examplex}[lemma]{Example}
\newtheorem{procedure}[lemma]{Procedure}

\newenvironment{exm}
  {\pushQED{\qed}\renewcommand{\qedsymbol}{$\diamond$}\examplex}
  {\popQED\endexamplex}
\newtheorem{question}[lemma]{Question}
\newtheorem{conjecture}[lemma]{Conjecture}
\newtheorem{Construction}[lemma]{Construction}
\newtheorem{algorithm}[lemma]{Algorithm}
%\theoremstyle{remark}
\newcommand{\define}[1]{{\bfseries\itshape #1}}
%%--MATH----------------------------------------------------------------------
\newcommand{\relphantom}[1]{\mathrel{\phantom{#1}}}
\newcommand{\abs}[1]{\ensuremath{\left| #1 \right|}}
%\newcommand{\norm}[1]{\ensuremath{\left\| #1 \right\|}}
\newcommand{\ideal}[1]{\ensuremath{\left\langle #1 \right\rangle}}
\renewcommand{\geq}{\geqslant}
\renewcommand{\leq}{\leqslant}
\newcommand{\coloneq}{\ensuremath{\mathrel{\mathop :}=}}
\newcommand{\transpose}{\ensuremath{\textsf{\upshape T}}} 
%% Blackboard bolds symbols
\renewcommand{\AA}{\ensuremath{\mathbb{A}}} 
\newcommand{\CC}{\ensuremath{\mathbb{C}}} 
\newcommand{\kk}{\ensuremath{\Bbbk}} 
\newcommand{\NN}{\ensuremath{\mathbb{N}}}
\newcommand{\PP}{\ensuremath{\mathbb{P}}} 
\newcommand{\QQ}{\ensuremath{\mathbb{Q}}} 
\newcommand{\RR}{\ensuremath{\mathbb{R}}} 
\newcommand{\EE}{\ensuremath{\mathbb{E}}} 

\renewcommand{\P}{\ensuremath{P}} 
\newcommand{\Q}{\ensuremath{Q}} 

%\renewcommand{\SS}{\ensuremath{\mathbb{S}}} 
\newcommand{\ZZ}{\ensuremath{\mathbb{Z}}} 
%% Script symbols
\newcommand{\sC}{\ensuremath{\mathcal{C}}}
\newcommand{\sE}{\ensuremath{\mathcal{E}}}
\newcommand{\sF}{\ensuremath{\mathcal{F}}}
\newcommand{\calF}{\ensuremath{\mathcal{F}}}
\newcommand{\calL}{\ensuremath{\mathcal{L}}}

\newcommand{\sH}{\ensuremath{\mathcal{H}}}
\newcommand{\sL}{\ensuremath{\mathcal{L}}}
\newcommand{\sO}{\ensuremath{\mathcal{O}}}
\newcommand{\diam}{\operatorname{diam}}
\newcommand{\Osh}{{\mathcal O}}
\newcommand{\Ish}{{\mathcal I}}
\newcommand{\calI}{{\mathcal I}}
\newcommand{\calH}{{\mathcal H}}

%% Operators
\DeclareMathOperator{\codim}{codim}
\DeclareMathOperator{\coker}{Coker}
\DeclareMathOperator{\conv}{conv}
\DeclareMathOperator{\Cox}{Cox}
\DeclareMathOperator{\depth}{depth}
\DeclareMathOperator{\gon}{gon}
\DeclareMathOperator{\Grass}{Gr}
\DeclareMathOperator{\green}{\textit{a}}
\DeclareMathOperator{\HH}{H}
\DeclareMathOperator{\HHH}{\mathbb{H}}
\DeclareMathOperator{\Hom}{Hom}
\DeclareMathOperator{\id}{id}
\DeclareMathOperator{\Image}{Im}
\DeclareMathOperator{\kept}{\lambda}
\DeclareMathOperator{\Ker}{Ker}
\DeclareMathOperator{\len}{\ell}
\DeclareMathOperator{\py}{py}
\DeclareMathOperator{\Pic}{Pic}
\DeclareMathOperator{\Proj}{Proj}
\DeclareMathOperator{\pr}{p}
\DeclareMathOperator{\qp}{qp}
\DeclareMathOperator{\qr}{qr}
\DeclareMathOperator{\rank}{rank}
\DeclareMathOperator{\reg}{reg}
\DeclareMathOperator{\Sos}{\Sigma}
\DeclareMathOperator{\Span}{Span}
\DeclareMathOperator{\Spec}{Spec}
\DeclareMathOperator{\sw}{sw}
\DeclareMathOperator{\Sym}{Sym}
\DeclareMathOperator{\topint}{int}
\DeclareMathOperator{\Tor}{Tor}
\DeclareMathOperator{\tw}{tw}
\DeclareMathOperator{\variety}{V}
%% Personalized comments
\newcommand\mv[1]{{\color{blue}(MV: #1)}}
\newcommand\mve[1]{{\color{red}(MV: #1)}}
\newcommand\jh[1]{{\color{red}(JH: #1)}}

\title{Ejercicios adicionales [E1]}

\begin{document}
\maketitle

\begin{enumerate}
\item {\it Medidas de probabilidad en conjuntos finitos.} Sea $\Omega$ un conjunto finito.
\begin{enumerate}
\item Demuestre que para toda $\sigma$-\'algebra $\calF$ sobre $\Omega$ existe una partición $V_1,\dots, V_m$ de $\Omega$ con $V_i\in \calF$ tal que todo conjunto en $\calF$ es unión disjunta de algunos $V_i$.
\item Demuestre que si $\alpha_1,\dots, \alpha_m$ son números reales no-negativos con $\sum \alpha_i=1$ entonces la asignación $\PP(V_j)=\alpha_j$ puede extenderse a una medida de probabilidad en $\PP: \calF\rightarrow [0,\infty)$ y que esta extensión es única.
\end{enumerate}

\item {\it Bernoulli trials con $n$ intentos.} Sea $n\in \mathbb{N}$ un entero positivo y sea $\Omega=\{0,1\}^n$ (es decir el conjunto de sucesiones de longitud $n$ con entradas en $\{0,1\}$).
Para $j\in \{1,\dots, n\}$ defina $A_j:=\{\omega\in \Omega: \omega_j=0\}$ y suponga que $p_j$ son números reales dados con $0\leq p_j\leq 1$.
\begin{enumerate}
\item Demuestre que existe una \'unica medida de probabilidad $\PP$ en los subconjuntos de $\Omega$ que cumple las siguientes condiciones:
\begin{enumerate}
\item Los eventos $A_1,\dots, A_n$ son $\PP$-independientes.
\item $\forall j\left(\PP(A_j)=p_j\right)$.
\end{enumerate}
\item Sea $B_k$ el evento correspondiente a las sucesiones con exáctamente $k$ ceros. Calcule $\PP(B_k)$ justificando su razonamiento.
\end{enumerate}

\item {\it Bernoulli trials en general.} Lanzamos una cantidad contable de monedas independientes, todas con probabilidad de cara $p$. Calcule la probabilidad de que la primera sucesión de tres símbolos idénticos consecutivos sea $HHH$ como función de $p$.

%\item {\it $2^{do}$ Borel-Cantelli.} Suponga que los $A_j$ son conjuntos $\PP$-independientes con $\PP(A_j)=\frac{1}{j+1}$. Demuestre que con probabilidad uno existen infinitos índices $n$ tales que $A_n$ ocurre pero $A_{n+1}$ no ocurre.(Sugerencia: Demuestre que los eventos $B_{2j}:=A_{2j}\setminus A_{2j+1}$ para $j=1,2,\dots$ son independientes)


\item {\it El problema de la ruina del apostador.} Un jugador comienza con $M$ unidades de dinero y juega partidas sucesivas. En cada partida gana una unidad con probabilidad $p$ ó pierde una unidad con probabilidad $q:=1-p$. Cuál es la probabilidad de que el jugador pierda todo su dinero?


\end{enumerate}


\end{enumerate}






\end{document}